\chapter{Vecteurs dans \texorpdfstring{$\R^2$ et $\R^3$}{R2 et R3}}
\label{workbook:chapter:09}
\section{Du point au vecteur}
\label{workbook:section:09.01}
\begin{exerciseblock}{Du point au vecteur}
\item Pour $A(-2,3)$ et $B(4,-1)$, calculer $\vect{AB}$ et $\vect{BA}$.
\item Donner deux segments orientés différents représentant le même vecteur $(3,2)$.
\item Expliquer la différence entre le point $(3,2)$ et le vecteur $(3,2)$, même si les coordonnées sont identiques.
\end{exerciseblock}
\section{Opérations sur les vecteurs}
\label{workbook:section:09.02}
\begin{exerciseblock}{Opérations sur les vecteurs}
\item Pour $u=(2,-3)$ et $v=(-5,4)$, calculer $u+v$, $u-v$, $3u$ et $-2v$.
\item Vérifier géométriquement et algébriquement la règle du parallélogramme pour $u=(3,1)$ et $v=(1,2)$.
\item Un marcheur effectue les déplacements $(4,1)$, $(-2,3)$ puis $(1,-5)$. Trouver le déplacement total.
\end{exerciseblock}
\section{Norme, distance et vecteur unitaire}
\label{workbook:section:09.03}
\begin{exerciseblock}{Norme, distance et vecteur unitaire}
\item Calculer la norme de $u=(6,-8)$ et le vecteur unitaire de même direction.
\item Calculer la distance entre $A(1,-2,3)$ et $B(4,2,-1)$.
\item Trouver tous les vecteurs de norme $5$ colinéaires à $(3,4)$.
\end{exerciseblock}
\section{Produit scalaire et angle}
\label{workbook:section:09.04}
\begin{exerciseblock}{Produit scalaire et angle}
\item Pour $u=(2,1)$ et $v=(-1,3)$, calculer $u\cdot v$ et l'angle entre eux.
\item Déterminer $k$ pour que $(k,2)$ soit orthogonal à $(3,-6)$.
\item Montrer que l'identité $\norm{u-v}^2=\norm u^2+\norm v^2-2u\cdot v$ contient la loi des cosinus.
\end{exerciseblock}
\section{Combinaisons linéaires}
\label{workbook:section:09.05}
\begin{exerciseblock}{Combinaisons linéaires}
\item Exprimer $(7,1)$ comme combinaison de $u=(1,1)$ et $v=(2,-1)$.
\item Déterminer si $(1,2,3)$ appartient au plan vectoriel engendré par $(1,0,1)$ et $(0,1,1)$.
\item Expliquer géométriquement pourquoi deux vecteurs non colinéaires de $\R^2$ engendrent tout le plan.
\end{exerciseblock}
\section{Interprétation géométrique : un déplacement libre}
\label{workbook:section:09.06}
\begin{exerciseblock}{Interprétation géométrique : un déplacement libre}
\item Translater le triangle de sommets $(0,0)$, $(2,0)$, $(1,3)$ par le vecteur $(-1,2)$.
\item Comparer deux flèches de même longueur et direction placées à des endroits différents : représentent-elles le même vecteur?
\item Une translation envoie $A$ sur $B$. Décrire l'image d'un point quelconque $P$ à l'aide de $\vect{AB}$.
\end{exerciseblock}
\section{Le produit scalaire comme longueur d'une ombre}
\label{workbook:section:09.07}
\begin{exerciseblock}{Le produit scalaire comme longueur d'une ombre}
\item Projeter $u=(4,3)$ sur la direction unitaire $e=(1,0)$, puis sur $f=(0,1)$.
\item Calculer la projection vectorielle de $u=(3,4)$ sur $v=(1,2)$.
\item Une force de $50$ N agit sur $8$ m avec un angle de $60^\circ$. Calculer le travail et interpréter la composante utile de la force.
\end{exerciseblock}
\section{Relier déplacements et coordonnées}
\label{workbook:section:09.08}
\begin{exerciseblock}{Relier déplacements et coordonnées}
\item Un drone part de $(2,-1,5)$ et subit les déplacements $(3,4,-2)$ puis $(-1,0,6)$. Trouver sa position finale.
\item Le milieu $M$ de $AB$ est caractérisé par $\vect{AM}=\frac12\vect{AB}$. En déduire la formule des coordonnées du milieu.
\item Déterminer le point $P$ divisant le segment de $A(1,2)$ à $B(7,5)$ dans le rapport $AP:PB=2:1$.
\end{exerciseblock}
\section*{Synthèse du chapitre}
\begin{synthesisbox}
\item Pour $A(1,0)$, $B(5,2)$ et $C(2,6)$, calculer les trois longueurs et l'angle en $A$, puis l'aire par déterminant ou trigonométrie.
\item Décomposer $u=(5,1)$ en une composante parallèle à $v=(2,2)$ et une composante orthogonale.
\item Montrer que les points $A(1,2,0)$, $B(3,1,4)$ et $C(5,0,8)$ sont alignés ou non en comparant les vecteurs.
\end{synthesisbox}
